% Nowe pozycje bibliografii — dopisać do środowiska thebibliography pracy
% jako kontynuację numeracji ([32]–[36]). Format zgodny z [1]–[31].
%
% Miejsca cytowań:
%   \cite{hardt2016}      -> 5.5b (definicja equalized odds, pierwsze użycie EOD)
%   \cite{bird2020}       -> 5.5b (narzędzie audytu fairlearn)
%   \cite{zadrozny2002}   -> 4.6  (kalibracja izotoniczna)
%   \cite{akiba2019}      -> 4.4.1 (strojenie Optuna) — patrz rozdzial4_instrukcje.md
%   \cite{grinsztajn2022} -> 5.2/5.3 (drzewa vs sieci na danych tabelarycznych)

\bibitem{hardt2016} Hardt M., Price E., Srebro N., ,,Equality of opportunity in supervised learning'' w \textit{Advances in Neural Information Processing Systems 29 (NeurIPS 2016)}, 2016, s. 3315--3323.
\bibitem{bird2020} Bird S., Dud\'{i}k M., Edgar R., Horn B., Lutz R., Milan V., Sameki M., Wallach H., Walker K., ,,Fairlearn: A toolkit for assessing and improving fairness in AI''. Microsoft, raport techniczny MSR-TR-2020-32, 2020.
\bibitem{zadrozny2002} Zadrozny B., Elkan C., ,,Transforming classifier scores into accurate multiclass probability estimates'' w \textit{Proceedings of the 8th ACM SIGKDD International Conference on Knowledge Discovery and Data Mining}, 2002, s. 694--699.
\bibitem{akiba2019} Akiba T., Sano S., Yanase T., Ohta T., Koyama M., ,,Optuna: A next-generation hyperparameter optimization framework'' w \textit{Proceedings of the 25th ACM SIGKDD International Conference on Knowledge Discovery and Data Mining}, 2019, s. 2623--2631.
\bibitem{grinsztajn2022} Grinsztajn L., Oyallon E., Varoquaux G., ,,Why do tree-based models still outperform deep learning on typical tabular data?'' w \textit{Advances in Neural Information Processing Systems 35 (NeurIPS 2022)}, 2022.
